% Online appendix. Compiled separately; does not count toward the 8,000-word limit.
% Tables come from tables/*.tex (make_tables.py and the gate scripts); figures from
% make_figures.py. Do not transcribe numbers by hand.
\documentclass[11pt,a4paper]{article}
\usepackage[utf8]{inputenc}
\usepackage[T1]{fontenc}
\usepackage{lmodern}
\usepackage[margin=1in]{geometry}
\usepackage{booktabs}
\usepackage{graphicx}
\usepackage{threeparttable}
\usepackage{microtype}
\PassOptionsToPackage{hyphens}{url}
\usepackage[hidelinks]{hyperref}
\renewcommand{\thetable}{A\arabic{table}}
\renewcommand{\thefigure}{A\arabic{figure}}
\title{Online Appendix\\[0.3em]\large Institutional Profiles Under a Common Prudential
Framework: Evidence from Brazilian Credit Cooperatives and Banks}
\author{Arthur Gomes Nery$^{a}$, Thiago de Oliveira Victorino$^{a,b}$ and
Rodrigo Lima Rangel$^{a}$\\[0.6em]
\small $^{a}$Organization of Brazilian Cooperatives (OCB System)\\
\small $^{b}$Department of Administration, University of Brasilia, Brazil}
\date{}
\begin{document}
\maketitle

\noindent All tables and figures are the authors' calculations from BACEN IF.data, on
the analysis sample of the main text unless a note says otherwise. Outcomes are defined
in Table~1 of the main text and the specifications in its Section~3.4. The material
follows the order of the results: the sample and the main estimates
(Tables~\ref{tab:a1} to~\ref{tab:a3}, Figure~\ref{fig:a1}), the dispersion of the
outcomes and the reporting-level split (Tables~\ref{tab:a4} and~\ref{tab:a5},
Figures~\ref{fig:a2} and~\ref{fig:a3}), the comparison groups (Tables~\ref{tab:a6}
and~\ref{tab:a7}), the network split (Table~\ref{tab:a8}), the sensitivity checks
(Tables~\ref{tab:a9} to~\ref{tab:a12}) and the de-cumulation of the income fields
(Figure~\ref{fig:a4}). The quantile regression coefficients are Table~5 of the main
text.

% ---------------------------------------------------------------- sample and main estimates
\begin{table}[htbp]\centering\begin{threeparttable}
\caption{Descriptive statistics}
\label{tab:a1}
\input{tables/t1_descriptives.tex}
\end{threeparttable}\end{table}

\begin{table}[htbp]\centering\begin{threeparttable}
\caption{Cooperative differential with confidence intervals}
\label{tab:a2}
\input{tables/t2_main_ci.tex}
\end{threeparttable}\end{table}

\begin{table}[htbp]\centering\begin{threeparttable}
\caption{Conditional mean and median}
\label{tab:a3}
\input{tables/t5_mean_vs_median.tex}
\end{threeparttable}\end{table}

\begin{figure}[htbp]
\centering
\includegraphics[width=0.8\textwidth]{figures/fig10_stability.png}
\caption{Cooperative coefficient by specification}
\label{fig:a1}
\par\smallskip{\footnotesize\raggedright In standard deviations of each outcome. Intervals are 95 percent, clustered by institution.\par}
\end{figure}

% ---------------------------------------------------------------- distribution and dispersion
\begin{table}[htbp]\centering\begin{threeparttable}
\caption{Dispersion at equal size}
\label{tab:a4}
\input{tables/t8_dispersion.tex}
\end{threeparttable}\end{table}

\begin{table}[htbp]\centering\begin{threeparttable}
\caption{Banks by reporting level}
\label{tab:a5}
\input{tables/t11_reporting_level.tex}
\end{threeparttable}\end{table}

\begin{figure}[htbp]
\centering
\includegraphics[width=0.8\textwidth]{figures/fig15_quantile_profiles.png}
\caption{Quantile coefficients relative to the conditional mean}
\label{fig:a2}
\par\smallskip{\footnotesize\raggedright Shaded: 20th to 80th percentile. The percentage beside each label is the average quantile coefficient over that range as a share of the conditional mean.\par}
\end{figure}

\begin{figure}[htbp]
\centering
\includegraphics[width=0.7\textwidth]{figures/fig16_dispersion.png}
\caption{Dispersion ratios at equal size}
\label{fig:a3}
\par\smallskip{\footnotesize\raggedright Cooperative over bank value, residualised on log assets and quarter.\par}
\end{figure}

% ---------------------------------------------------------------- comparison groups
\begin{table}[htbp]\centering\begin{threeparttable}
\caption{Cooperative differential by comparison group}
\label{tab:a6}
\input{tables/t3_peer_ladder.tex}
\end{threeparttable}\end{table}

\begin{table}[htbp]\centering\begin{threeparttable}
\caption{Provisioning by capital tercile}
\label{tab:a7}
\input{tables/t12_provisioning_by_capital.tex}
\end{threeparttable}\end{table}

% ---------------------------------------------------------------- networks
\begin{table}[htbp]\centering\begin{threeparttable}
\caption{Cooperative differential by system}
\label{tab:a8}
\input{tables/t4_systems.tex}
\end{threeparttable}\end{table}

% ---------------------------------------------------------------- sensitivity
\begin{table}[htbp]\centering\begin{threeparttable}
\caption{Excluding S1 and S2 banks}
\label{tab:a9}
\input{tables/t6_segment_restricted.tex}
\end{threeparttable}\end{table}

\begin{table}[htbp]\centering\begin{threeparttable}
\caption{Before and after the 2022 capital rules}
\label{tab:a10}
\input{tables/t9_subperiod.tex}
\end{threeparttable}\end{table}

\begin{table}[htbp]\centering\begin{threeparttable}
\caption{Return on assets before and after taxes}
\label{tab:a11}
\input{tables/t13_pretax.tex}
\end{threeparttable}\end{table}

\begin{table}[htbp]\centering\begin{threeparttable}
\caption{Stability outcomes before and after taxes}
\label{tab:a12}
\input{tables/t15_pretax_stability.tex}
\end{threeparttable}\end{table}

% ---------------------------------------------------------------- measurement
\begin{figure}[htbp]
\centering
\includegraphics[width=0.75\textwidth]{figures/fig17_decumulation.png}
\caption{Return on assets before and after de-cumulation}
\label{fig:a4}
\par\smallskip{\footnotesize\raggedright Median across the analysis sample, annualised.\par}
\end{figure}

\end{document}
